\documentclass[10pt]{article}
\usepackage[margin=0.9in]{geometry}
\usepackage{amsmath,microtype}
\title{Guide: What Sampled Contraction Would Prove}
\author{Collatz proof project}
\date{September 5, 2026}
\begin{document}
\maketitle
\raggedright
\section*{The exact result}
Lean now proves that the following two claims are equivalent: every
positive Collatz orbit is bounded; and every positive seed two modulo
nine eventually has a coefficient below one at a later visit to that
same progression. Both claims remain unproved.

\section*{Why restricting the observations works}
Suppose the coefficient never contracts at any of those visits.
Between consecutive visits, at most six even steps occur. Even steps
halve the coefficient and odd steps increase it, so at every intervening
time it remains at least one sixty-fourth. Our previous exact correction
theory then shows that such a positive seed must have an unbounded orbit.

Conversely, on any hypothetical unbounded orbit the inverse coefficient
tends to zero. Look only at visits to two modulo nine and choose a visit
where this inverse coefficient is largest. Relative to that starting
point, every later sampled coefficient is at least one. This supplies
both the required residue and sampled noncontraction in the same tail.

\section*{What remains unresolved}
The selected tail need not have coefficients at least one at every time;
the proved intervening bound is one sixty-fourth. The distinction matters
when applying earlier lemmas that require full noncontraction.

We still need a proof that every seed in the progression eventually
contracts at a visit. The equivalence identifies that as the full
nondivergence problem. It does not eliminate nontrivial cycles.
The finite experiment about descent at the first sampled contraction
remains a separate, unproved candidate. No proof of Collatz is claimed.

\section*{Verification}
From \texttt{lean/}, run
\texttt{lake build Collatz.SampledContraction}.
The repository's selected theorem audit checks the four new declarations
and their axiom dependencies. The proof uses existing exact theorems,
without importing an experimental result as an assumption.
\end{document}
